\section{Introduction}
\done\todo{Define the problem: lack of real-world examples to teach and learn software security. This project aims to curate a set of software vulnerability examples.}
\done\todo{Discuss the uniqueness of the educational needs: (a) precision (correct examples), (b) length and complexity; (c) prerequisite knowledge.}
\done\todo{Discuss the overarching method to do it: (a) retrieve patches from NVD; (b) determine type of software vulnerability; (c) determine prerequisite; (d) score difficulties -- (c) and (d) are more open questions.}
Software security has become a pillar of modern cybersecurity and software engineering education due to the increasing frequency and impact of real-world vulnerabilities and their exploits. \attention{false}Hence, a greater emphasis is placed upon instructors to provide proactive measures to prevent vulnerabilities by preparing next-generation software developers. Prior work has shown that automated scaffolding has led to improved student reflection and understanding in security contexts~\cite{chen2019using}. In order to do so efficiently, \attention{false}{we argue that} students and instructors require real-world examples to learn and teach from.  \attention{false}{Research suggests that} students who work with real-world examples of software vulnerabilities report a greater appreciation for the opportunity to view the examples in detail and that there is an overall gain in self-efficacy after using those examples~\cite{fie2024dlcd}. Yet, modern undergraduate cybersecurity and software engineering programs still lack a systematically curated collection of vulnerabilities to effectively teach prevention measures.

The \attention{false}{ACM/IEEE/AAAI} Computer Science Curricula 2023 (CS2023) revises and expands on the \attention{false}{ACM/IEEE} CS2013 guidelines~\cite{cscurriculum2023}, revealing a contradiction: although CS2023 explicitly emphasizes the necessity of a security mindset \cite[page 256]{cscurriculum2023}, its Secure Coding section provides no accompanying database or set of vulnerabilities for instructional use \cite[page 259]{cscurriculum2023}. As such, instructors lack a purpose-built and robust set of patched and thoroughly documented vulnerabilities, and may need to resort to outdated or overly simplified examples that are unable to capture the severity of vulnerabilities and the importance of embedding security into code. Consequently, students may graduate with theoretical or simplified knowledge on common vulnerabilities, but they will not have the contextual understanding to properly address these issues in a continually evolving field.

The National Vulnerability Database (NVD)\attention{false}{~\cite{nvd}} currently holds hundreds of thousands of detailed Common Vulnerability and Exposure (CVE) entries that are linked to real-world projects, some of which contain real patches and solutions\attention{false}{~\cite{nvddashboard}}. While many CVE entries contain Common Weakness Enumeration (CWE) classifications for vulnerability types, a portion of entries do not. \attention{false}{In addition,} these entries are not organized by instructional complexity or by updated didactically relevant patches; some lack patches entirely or reference outdated, extensive patches that are unsuitable for some programmers \cite{nvidcweslice}. Moreover, patches that have been published are rarely updated to ensure accuracy over time. \deletetext{true}{Lastly, as patches are derived exclusively from open-source projects, no high-quality commercial examples exist. }Ultimately, this leads to a dual issue of missing CWE classifications and the lack of educational organization, which creates barriers for instructors and students who seek proper examples to use in their education.

\deletetext{true}{Despite the NVD's current repository of thousands of real-world CVEs, it is incapable of meeting the needs of instructors and students in the pursuit of a greater security mindset. Entries contain no organization by educational complexity, outdated or missing patches, and exclusive reliance on open-source repositories. }Our proposal outlines a methodology to fill these shortcomings by filtering for updated patches, classifying vulnerability types, denoting prerequisite concepts, and assigning didactic difficulty levels.